\section{Training Procedure}
\label{sec:training}

\subsection{Training Data Generation}

Training data is generated by running PCGCv2 on the 8iVFB dataset at rates
r2, r4, and r6, decoding to produce $\hat{\pc}_t$ for each frame $t$.
Three of the four sequences are used for training: \textit{longdress},
\textit{loot}, and \textit{soldier}.
The fourth sequence (\textit{redandblack}) is reserved as the held-out
validation sequence throughout all experiments.

\paragraph{Frame sampling.}
Adjacent frames of the same sequence are highly correlated (a subject's pose
changes slowly at 10 fps).
Using all 300 frames per sequence would provide little additional training
signal while increasing dataset redundancy and storage requirements.
We sample every 12th frame (frame stride 12), yielding approximately 25 frames
per sequence per rate, for a total of $3 \times 3 \times 25 = 225$ training
clouds.

\paragraph{Patch extraction.}
Each cloud is divided into overlapping $64^3$ voxel patches with stride 32
(50\% overlap).
A typical 8iVFB cloud at $G = 1024$ with 500\,000 points yields approximately
500--2000 patches depending on the sequence and compression rate.
Empty patches (those containing fewer than 10 decoded points) are discarded.

\paragraph{Data augmentation.}
The cubic symmetry group has 48 elements~\cite{winkels2018gcnn,cohen2016equivariant}: all permutations ($3! = 6$) and
sign flips ($2^3 = 8$) of the three axes.
Each patch is randomly assigned one of the 48 discrete orientations at
loading time, providing orientation invariance without requiring a learned
equivariant backbone.
The augmentation is label-preserving: the same transformation is applied to
both the decoded patch and the original crop.

\subsection{Optimiser and Learning Rate Schedule}

\begin{description}
  \item[Optimiser:] AdamW~\cite{loshchilov2017decoupled} with learning rate
    $\eta = 10^{-3}$ and weight decay $\lambda = 10^{-4}$.
    The $\beta_1 = 0.9$, $\beta_2 = 0.999$ default Adam moment parameters
    are used.

  \item[Warm-up:] Linear increase from $\eta_0 = 10^{-5}$ to $10^{-3}$ over
    3 epochs (approximately 3\,000 gradient steps).
    Warm-up prevents large gradient steps before batch normalisation statistics
    have converged.

  \item[Cosine annealing:] After warm-up, the learning rate follows:
    $\eta_t = \eta_{\min} + \frac{1}{2}(10^{-3} - \eta_{\min})(1 + \cos(\pi t / 27))$
    for $t = 3, \ldots, 30$ epochs, with $\eta_{\min} = 10^{-6}$.
\end{description}

\subsection{Hardware and Runtime}

Training runs on a single NVIDIA RTX 3090 GPU (24 GB VRAM).
The batch size is 8 patches per mini-batch (approximately 3 patches per rate,
varying due to random sampling).
Each epoch requires approximately 2 hours for the $\approx$225 training clouds.
Total training time for 30 epochs is approximately 60 hours.

\subsection{Checkpoint Selection}

Model checkpoints are saved after each epoch.
The best checkpoint is selected by validation loss on \textit{redandblack}
(which is excluded from training).
For the final submitted model (v10-nocurv), the best checkpoint is at epoch 24
out of 30.
The epoch-24 checkpoint is used in all reported evaluations.

\subsection{Training Stability}

The training loss converges smoothly without significant instability.
Figure~\ref{fig:training-curve} shows the training and validation loss curves
for the v10-nocurv model.
The validation loss (on \textit{redandblack}) tracks the training loss closely
throughout, with no significant overfitting.
This is expected given the data augmentation (48 orientations effectively
multiplying the training set size) and the relatively small model size for
the parameter count.

\begin{figure}[t]
  \centering
  \includegraphics[width=0.82\textwidth]{training_curve.png}
  \caption{%
    Training and validation Chamfer Distance loss curves for the v10-nocurv model.
    The warm-up phase (shaded, epochs 1--3) increases the learning rate linearly
    from $10^{-5}$ to $10^{-3}$; cosine annealing follows for epochs 4--30.
    The best checkpoint (vertical dashed line) is at epoch 24, selected by
    minimum validation loss on the held-out \textit{redandblack} sequence.
    The close tracking of training and validation loss indicates no significant
    overfitting.
  }
  \label{fig:training-curve}
\end{figure}

\subsection{Full Model Summary}

Table~\ref{tab:model-summary} summarises the architecture and training
configuration of the final submitted model.

% ============================================================
